\documentclass[11pt]{article}

\usepackage[margin=1in]{geometry}
\usepackage[T1]{fontenc}
\usepackage[utf8]{inputenc}
\usepackage{amsmath,amssymb}
\usepackage{booktabs}
\usepackage{longtable}
\usepackage{array}
\usepackage[table]{xcolor}
\usepackage{adjustbox}
\usepackage{pdflscape}
\usepackage{parskip}
\usepackage{enumitem}
\usepackage{hyperref}
\hypersetup{colorlinks=true, linkcolor=blue, citecolor=blue, urlcolor=blue}

% --- external cross-references into the main manuscript ---
% Compile jmlr_paper_main_3107919.tex FIRST (to produce its .aux file),
% then compile this document, for \ref{...} below to resolve to the main
% paper's actual section/table numbers rather than "??".
\usepackage{xr}
\externaldocument{jmlr_paper_main_3107919}

\newcommand{\rtwo}{$R^2$}
\newcommand{\open}{\textcolor{red!60!black}{\textbf{[OPEN]}}}
\newcommand{\closed}{\textcolor{green!40!black}{\textbf{[CLOSED]}}}
\newcommand{\flag}{\textcolor{orange!80!black}{\textbf{[FLAG]}}}
\newcommand{\gap}{\textcolor{purple!60!black}{\textbf{[DATA GAP]}}}
\newcommand{\Rsq}{R^{2}}

\definecolor{match}{HTML}{C6EFCE}
\definecolor{minor}{HTML}{FFEB9C}
\definecolor{major}{HTML}{FFC7CE}
\newcommand{\legenditem}[2]{\colorbox{#1}{\phantom{XX}}~#2}

\newenvironment{appliesto}[1]
  {\begin{quote}\small\itshape\noindent Applies to main manuscript
    (\texttt{jmlr\_paper\_main\_3107919.tex}): #1\end{quote}}
  {}

\title{\bf Consolidated Reproducibility \& Audit Supplement\\
\large Cross-Referenced to \emph{HypatiaX: A Hybrid Symbolic--Neural
Framework for Extrapolation-Reliable Analytical Discovery}}
\author{}
\date{Compiled 2026-07-30}

\begin{document}
\maketitle
\thispagestyle{empty}

\begin{abstract}
This document consolidates six previously separate audit and
reproducibility notes into a single cross-referenced supplement. Each
section below states, up front, which section(s) and table(s) of the
main manuscript it corrects, extends, or provides evidence for. Where
the main manuscript already reflects a finding from one of these notes,
that is stated explicitly rather than left implicit. Sections are
ordered to follow the manuscript's own section order, ending with a
cross-cutting edit tracker that indexes every item back to the sections
above.

\medskip
\noindent\textbf{Contents:}
\begin{enumerate}[nosep]
\item Runtime Table 1 --- full by-seed / by-split reconstruction
  (\S\S~1--1, corrects \ref{sec:timing})
\item Core-15 Ablation (``Table~6'') --- real-data vs.\ published
  comparison (\ref{sec:ablation_core15})
\item Feynman-30 Extrapolation Benchmark, Section~10.7 update
  (\ref{sec:feynman30})
\item Reproducibility note: unpinned LLM model string, Nguyen-12 pipeline
  (\ref{sec:nguyen12})
\item CI dashboard status, Item 10.1 (\ref{sec:process_items})
\item Consolidated manuscript-edit tracker (cross-cutting)
\end{enumerate}
\end{abstract}

\clearpage

% =======================================================================
\section{Runtime Table 1: Full By-Seed / By-Split Reconstruction}
\label{supp:runtime}
% =======================================================================

\begin{appliesto}
\S\ref{sec:timing} (``Runtime Analysis''), Table~\ref{tab:timing_full},
Table~\ref{tab:timing_llm_routed_full}, and the top-of-paper runtime note.
\textbf{Status: already incorporated} --- the main manuscript's current
Table~\ref{tab:timing_full} is the direct output of this note. Reproduced
here so the full by-seed evidence trail travels with the supplement
rather than only living in the main file's edit history.
\end{appliesto}

\subsection{Table 1 (full) --- Runtime by Seed and by Split}

Regenerated directly from the raw per-task result files by
\texttt{generate\_table1.py}, with no hand-edited cells. Rows are grouped
by split (aggressive 40/60 = v3c, PCA-directed 40/60), with one row per
seed (42, 99, 123, 777, 2024) followed by the pooled 5-seed / 370-task
total for that split (shaded). ``Speedup'' compares the neural MLP arm to
the hybrid arm; values labeled SLOWER mean hybrid took \emph{longer}, not
less.

\begin{table}[h]
\centering
\rowcolors{2}{white}{white}
\resizebox{\linewidth}{!}{%
\begin{tabular}{lrrrrrr}
\toprule
Run & $n$ & Pure LLM (s) & Neural MLP (s) & Hybrid, all (s) & LLM-routed & Speedup, all \\
    &     & mean/median  & mean/median    & mean/median      & count      & (mean) \\
\midrule
v3c seed42  & 74 & 8.85 / 7.75 & 0.36 / 0.36 & 2.54 / 1.49 & 69/74 & 7.05x SLOWER \\
v3c seed99  & 74 & 8.56 / 7.62 & 0.37 / 0.37 & 2.52 / 1.48 & 69/74 & 6.83x SLOWER \\
v3c seed123 & 74 & 8.46 / 7.30 & 0.36 / 0.36 & 2.53 / 1.58 & 68/74 & 7.06x SLOWER \\
v3c seed777 & 74 & 8.66 / 7.67 & 0.25 / 0.25 & 2.32 / 1.42 & 69/74 & 9.33x SLOWER \\
v3c seed2024 & 74 & 8.35 / 7.75 & 0.37 / 0.36 & 2.31 / 1.38 & 69/74 & 6.33x SLOWER \\
\rowcolor{blue!8}
\textbf{v3c ALL SEEDS (pooled)} & \textbf{370} & \textbf{8.57 / 7.65} & \textbf{0.34 / 0.36} & \textbf{2.44 / 1.46} & \textbf{344/370} & \textbf{7.18x SLOWER} \\
\midrule
PCA seed42  & 74 & 8.44 / 7.37 & 0.36 / 0.36 & 2.18 / 1.40 & 68/74 & 6.11x SLOWER \\
PCA seed99  & 74 & 2.97 / 0.19 & 0.35 / 0.35 & 1.16 / 0.91 & 26/74 & 3.28x SLOWER \\
PCA seed123 & 74 & 0.21 / 0.17 & 0.36 / 0.35 & 0.90 / 0.88 & 0/74  & 2.52x SLOWER \\
PCA seed777 & 74 & 0.19 / 0.17 & 0.36 / 0.35 & 0.91 / 0.88 & 0/74  & 2.55x SLOWER \\
PCA seed2024 & 74 & 0.24 / 0.17 & 0.36 / 0.35 & 0.90 / 0.88 & 0/74 & 2.54x SLOWER \\
\rowcolor{blue!8}
\textbf{PCA ALL SEEDS (pooled)} & \textbf{370} & \textbf{2.41 / 0.18} & \textbf{0.36 / 0.35} & \textbf{1.21 / 0.89} & \textbf{94/370} & \textbf{3.40x SLOWER} \\
\bottomrule
\end{tabular}}
\caption{Wall-clock time per task (seconds), by seed and by split, plus
pooled 5-seed totals. No timeouts were recorded on any arm, any seed,
either split. This is the same table as the main manuscript's
Table~\ref{tab:timing_full}.}
\label{supp:tab:timing_full}
\end{table}

\subsection{Reading the by-seed breakdown}

Within the v3c (aggressive) split, hybrid is slower than the neural
network on \emph{every single seed} --- the slowdown ranges from about
$6.3\times$ (seed 2024) to about $9.3\times$ (seed 777) on mean time, with
no seed showing hybrid faster. This is not a pooling artifact: it is true
seed-by-seed.

Within the PCA-directed split, the picture splits into two regimes
depending on the seed. Seeds 42 and 99 route a meaningful share of tasks
to the LLM (68/74 and 26/74 respectively) and show correspondingly larger
hybrid slowdowns ($6.1\times$ and $3.3\times$ on mean). Seeds 123, 777, and
2024 route \emph{zero} tasks to the LLM under this router
($n_{\text{llm\_routed}} = 0$), yet hybrid is still slower than the NN on
those seeds too (about $2.5\times$ on mean) --- meaning the hybrid wrapper
itself carries a fixed overhead beyond just the LLM call, since with zero
LLM calls hybrid should be timing-equivalent to the NN and it is not.

Comparing the two pooled split totals: PCA-directed routing produces a
smaller hybrid slowdown overall ($3.40\times$ mean) than the aggressive
split ($7.18\times$ mean), consistent with PCA calling the LLM far less
often (94/370 vs.\ 344/370 tasks). But neither split shows hybrid
outperforming the neural network at any level of aggregation --- by seed,
pooled, or restricted to LLM-routed tasks only.

\subsection{What the experiment is testing}

The benchmark compares three ways of answering DeFi-related tasks: a pure
LLM call, a lightweight neural network (MLP) classifier, and a ``hybrid''
system that routes each task to either the NN or the LLM depending on a
decision rule. Two routing strategies are tested --- an aggressive 40/60
split (v3c) and a PCA-directed 40/60 split --- each run across five seeds
(42, 99, 123, 777, 2024) with 74 tasks per seed, for 370 tasks total per
split. The stated goal of the hybrid design is to capture most of the
speed of the NN while keeping the LLM's judgment available for harder
cases --- i.e., faster than pure LLM, and not much slower than pure NN.

\subsection{Where this diverges from an earlier draft}

An earlier draft of the manuscript reported the opposite pattern: hybrid
faster than NN, by a claimed $1.73\times$ (mean) / $1.64\times$ (median),
with NN itself reported around 3.0\,s mean / 2.7\,s median and hybrid
around 6.8\,s mean / 1.7\,s median. Those absolute NN numbers are also
$8$--$9\times$ larger than what appears in any raw file available here
(raw NN times cluster around 0.34--0.42\,s). That gap is large enough to
look less like a computation error and more like the reported numbers
came from a different run entirely --- a heavier or older NN
implementation, a different task set, or a stale draft of the table that
was never updated after the current scripts were finalized. An even
earlier draft additionally claimed that seed files other than seed~42
were truncated to 5 of 74 tasks; that claim is also retracted --- all ten
raw files (five seeds $\times$ two splits) contain the full 74 tasks
each, correctly seed-labeled with no cross-contamination.

\subsection{Net effect}

The hybrid design does not deliver a speedup over the neural network in
this data; it is slower, sometimes substantially so, and the size of
that slowdown tracks cleanly with how often each router calls the LLM.
The PCA router being meaningfully better than the aggressive router at
limiting this overhead is a real and defensible finding worth reporting
on its own terms --- ``PCA-directed routing reduces but does not
eliminate hybrid's overhead relative to NN-only'' --- but it is a
different, more modest claim than the $1.7\times$ speedup in the
withdrawn earlier draft.

\clearpage
% =======================================================================
\section{Core-15 Ablation (``Table 6''): Real-Data vs.\ Published Comparison}
\label{supp:table6}
% =======================================================================

\begin{appliesto}
\S\ref{sec:ablation_core15} (``Ablation: PySR-only vs.\ HypatiaX
(Core-15)''), Table~\ref{tab:llm_ablation}. \textbf{Status: already
incorporated} --- the main manuscript's Table~\ref{tab:llm_ablation} text
states it ``has been rebuilt from the underlying \texttt{\_merged.json}
data \ldots replacing the withdrawn version.'' This section is the
cell-by-cell discrepancy evidence behind that rebuild, preserved here for
audit trail. If Table~\ref{tab:llm_ablation} is ever re-edited, it should
be checked against the \emph{Real} columns below, not against the
\emph{Pub} columns, which are the withdrawn transcription.
\end{appliesto}

\begin{landscape}

\subsection*{Table 6 (as Published) vs.\ Real \texttt{\_merged.json} Values --- Core-15 Ablation}

\textbf{Legend:} \quad
\legenditem{match}{published matches real data (within rounding)} \qquad
\legenditem{minor}{notable numeric discrepancy} \qquad
\legenditem{major}{major discrepancy (sign flip, order-of-magnitude gap, or fabricated/missing value)}

\bigskip
\noindent Coloring is applied to the \emph{Real} columns, relative to the
corresponding \emph{Pub}(lished) value in the same row/metric. P~=~PySR-only;
H~=~HypatiaX.

\bigskip

\begin{adjustbox}{width=\linewidth,center}
\begin{tabular}{l l *{4}{r} *{4}{r} *{4}{r} *{4}{r}}
\toprule
& & \multicolumn{4}{c}{Train $R^2$} & \multicolumn{4}{c}{Near (1.2$\times$) $R^2$} & \multicolumn{4}{c}{Medium $R^2$} & \multicolumn{4}{c}{Far (5$\times$) $R^2$} \\
\cmidrule(lr){3-6} \cmidrule(lr){7-10} \cmidrule(lr){11-14} \cmidrule(lr){15-18}
Equation & Domain & Pub P & Pub H & Real P & Real H & Pub P & Pub H & Real P & Real H & Pub P & Pub H & Real P & Real H & Pub P & Pub H & Real P & Real H \\
\midrule
\textbf{Allometric Scaling} & Biology & 0.9977 & 0.9973 & \cellcolor{match}0.9977 & \cellcolor{match}0.9979 & 0.9996 & 0.9509 & \cellcolor{match}1.0000 & \cellcolor{minor}0.9824 & 1.0000 & 0.8602 & \cellcolor{match}1.0000 & \cellcolor{minor}0.9502 & 0.9996 & -2.1139 & \cellcolor{match}1.0000 & \cellcolor{major}-0.0644 \\
\textbf{Arrhenius} & Chemistry & 0.9896 & 0.9971 & \cellcolor{match}0.9974 & \cellcolor{match}0.9979 & -0.9783 & -0.4012 & \cellcolor{major}0.9779 & \cellcolor{major}0.9977 & -0.6766 & -0.6624 & \cellcolor{major}0.4636 & \cellcolor{major}0.9957 & -12.5549 & -12.5553 & \cellcolor{match}-12.5552 & \cellcolor{major}0.9028 \\
\textbf{Constant Product} & DeFi AMM & 0.9982 & 0.9982 & \cellcolor{match}0.9982 & \cellcolor{match}0.9985 & 0.9996 & 0.9996 & \cellcolor{match}0.9996 & \cellcolor{major}-18{,}071.5 & 1.0000 & 1.0000 & \cellcolor{match}1.0000 & \cellcolor{major}-5.4$\times10^{34}$ & 0.9996 & 0.9996 & \cellcolor{match}0.9996 & \cellcolor{major}-1.9$\times10^{46}$ \\
\textbf{Gravitational Force} & Physics & 0.9146 & 0.9544 & \cellcolor{minor}0.9976 & \cellcolor{minor}0.9889 & -4.2880 & -2.6752 & \cellcolor{major}0.9973 & \cellcolor{major}$-\infty$ & -0.0260 & -0.0016 & \cellcolor{major}0.9998 & \cellcolor{major}$-\infty$ & -9.0418 & -7.6360 & \cellcolor{major}0.9973 & \cellcolor{major}$-\infty$ \\
\textbf{Henderson-Hasselbalch} & Chemistry & 0.9123 & 0.9338 & \cellcolor{minor}0.9339 & \cellcolor{minor}0.9986 & 0.2137 & 0.2172 & \cellcolor{major}0.8004 & \cellcolor{major}\textit{null} & 0.9633 & -3.6019 & \cellcolor{minor}0.9850 & \cellcolor{major}\textit{null} & 0.2137 & -4.9142 & \cellcolor{major}0.8004 & \cellcolor{major}\textit{null} \\
\textbf{Ideal Gas Law} & Physics & 0.9976 & 0.9976 & \cellcolor{match}0.9976 & \cellcolor{match}0.9980 & 0.9999 & 0.9999 & \cellcolor{match}1.0000 & \cellcolor{match}0.9988 & 1.0000 & 1.0000 & \cellcolor{match}1.0000 & \cellcolor{match}0.9999 & 0.9999 & 0.9999 & \cellcolor{match}1.0000 & \cellcolor{match}0.9993 \\
\textbf{Impermanent Loss} & DeFi AMM & 0.9975 & 0.9975 & \cellcolor{match}0.9971 & \cellcolor{match}0.9977 & 0.9121 & 0.9113 & \cellcolor{minor}0.7703 & \cellcolor{minor}0.9639 & -0.3063 & -0.3091 & \cellcolor{major}-2.7832 & \cellcolor{major}-80.4327 & -62.4026 & -62.5166 & \cellcolor{major}-157.0776 & \cellcolor{major}-8{,}259.3 \\
\textbf{Kinetic Energy} & Physics & 0.9968 & 0.9968 & \cellcolor{match}0.9968 & \cellcolor{match}0.9968 & 1.0000 & 1.0000 & \cellcolor{match}1.0000 & \cellcolor{match}1.0000 & 1.0000 & 1.0000 & \cellcolor{match}1.0000 & \cellcolor{match}1.0000 & 1.0000 & 1.0000 & \cellcolor{match}1.0000 & \cellcolor{match}1.0000 \\
\textbf{Liquidation Price} & DeFi Risk & 0.9974 & 0.9974 & \cellcolor{match}0.9974 & \cellcolor{match}0.9977 & 0.9999 & 1.0000 & \cellcolor{match}1.0000 & \cellcolor{match}0.9998 & 1.0000 & 1.0000 & \cellcolor{match}1.0000 & \cellcolor{match}0.9999 & 1.0000 & 1.0000 & \cellcolor{match}1.0000 & \cellcolor{match}0.9993 \\
\textbf{Logistic Growth} & Biology & 0.9974 & 0.9975 & \cellcolor{match}0.9975 & \cellcolor{match}0.9976 & 0.9795 & 0.9999 & \cellcolor{minor}0.9997 & \cellcolor{major}\textit{null} & 0.9947 & 1.0000 & \cellcolor{match}0.9998 & \cellcolor{major}\textit{null} & 0.9934 & 0.9999 & \cellcolor{match}0.9996 & \cellcolor{major}\textit{null} \\
\textbf{Michaelis-Menten} & Biology & 0.9948 & 0.9968 & \cellcolor{match}0.9975 & \cellcolor{match}0.9979 & -68.5896 & -0.0979 & \cellcolor{major}0.9890 & \cellcolor{major}\textit{null} & -368.7928 & -2.4717 & \cellcolor{major}0.9951 & \cellcolor{major}\textit{null} & -83{,}899.5 & -634.5989 & \cellcolor{major}0.6471 & \cellcolor{major}\textit{null} \\
\textbf{Portfolio Std Dev} & DeFi Risk & 0.9504 & 0.9975 & \cellcolor{minor}0.9975 & \cellcolor{match}0.9975 & 0.8865 & 1.0000 & \cellcolor{minor}1.0000 & \cellcolor{match}1.0000 & 0.9493 & 1.0000 & \cellcolor{minor}1.0000 & \cellcolor{match}1.0000 & -118.4482 & 1.0000 & \cellcolor{major}1.0000 & \cellcolor{match}1.0000 \\
\textbf{Price Impact} & DeFi AMM & 0.9976 & 0.9976 & \cellcolor{match}0.9976 & \cellcolor{match}0.9976 & 1.0000 & 1.0000 & \cellcolor{match}1.0000 & \cellcolor{match}1.0000 & 1.0000 & 1.0000 & \cellcolor{match}1.0000 & \cellcolor{match}1.0000 & 1.0000 & 1.0000 & \cellcolor{match}1.0000 & \cellcolor{match}1.0000 \\
\textbf{Rate Law} & Chemistry & 0.9977 & 0.9977 & \cellcolor{match}0.9977 & \cellcolor{match}0.9981 & 1.0000 & 0.9999 & \cellcolor{match}1.0000 & \cellcolor{major}\textit{null} & 1.0000 & 1.0000 & \cellcolor{match}1.0000 & \cellcolor{major}\textit{null} & 1.0000 & 0.9999 & \cellcolor{match}1.0000 & \cellcolor{major}\textit{null} \\
\textbf{Value at Risk} & DeFi Risk & 0.9979 & 0.9979 & \cellcolor{match}0.9979 & \cellcolor{match}0.9979 & 0.9999 & 0.9999 & \cellcolor{match}0.9999 & \cellcolor{match}1.0000 & 1.0000 & 1.0000 & \cellcolor{match}1.0000 & \cellcolor{match}1.0000 & 0.9999 & 0.9999 & \cellcolor{match}0.9999 & \cellcolor{match}1.0000 \\
\bottomrule
\end{tabular}
\end{adjustbox}

\bigskip

{\small
\textbf{Sources.} \emph{Pub} columns transcribed verbatim from the
withdrawn draft of Table~6 (formerly Supplementary~A \S10.6).
\emph{Real} columns computed directly from \texttt{\_merged.json}
(\texttt{hypatia} = H, \texttt{pysr\_only} = P fields), read
programmatically with no intermediate transcription step.
}

\subsection*{Notes}

\begin{enumerate}
\item \textbf{Arrhenius} --- Published reports $H_{\text{far}} = -12.5553$, essentially copying P's real failure value. Real $H_{\text{far}} = +0.9028$ (a success). A discussion built on the withdrawn ``both PySR-only and HypatiaX collapse catastrophically \ldots identical failure mechanism'' framing does not match the underlying data for H.

\item \textbf{Michaelis-Menten} --- The withdrawn draft's headline claims $P_{\text{far}} = -83{,}900$ and $H_{\text{far}} = -635$ (a ``132$\times$ reduction''). Real data: $P_{\text{far}} = +0.647$ (a clear success, not catastrophic), and $H$ is \textit{null} at near/med/far (evaluation never produced a value). Neither headline number appears in the raw data.

\item \textbf{Rate Law \& Logistic Growth} --- Real $H$ is \textit{null} across near/med/far for both (evaluation failed). The withdrawn draft reports near-perfect $H$ values ($\approx 0.9999$--$1.0000$) that correspond to no recorded run, and are suspiciously close to the real $P$ values.

\item \textbf{Gravitational Force} --- Real $H$ is $-\infty$ across near/med/far (a genuine crash); real $P$ is $\approx 0.997$--$0.9998$ (near-perfect). The withdrawn draft reports finite, moderate-negative numbers for \emph{both} columns --- neither the real $P$ success nor the real $H$ blow-up appear.

\item \textbf{Constant Product} --- Real $H$ collapses to $-5.4\times10^{34}$ (medium) and $-1.9\times10^{46}$ (far), but the withdrawn draft shows $H$ tracking $P$ almost exactly at 0.9996 throughout.

\item \textbf{The P (PySR-only) column is also affected} --- Several withdrawn $P$ values for near/medium ranges (Arrhenius near/med, Henderson-Hasselbalch near, Michaelis-Menten near/med) are substantially different from the real $P$ values, and in most cases the withdrawn number makes PySR-only look \emph{worse} than it really performed. Combined with the H-column errors, this widened HypatiaX's apparent margin from both directions at once in the withdrawn draft, rather than from a single column's errors.
\end{enumerate}

\end{landscape}

\clearpage
% =======================================================================
\section{Feynman-30 Extrapolation Benchmark: Section 10.7 Update}
\label{supp:feynman30}
% =======================================================================

\begin{appliesto}
\S\ref{sec:feynman30} (``Feynman Extrapolation Benchmark''), \S
\ref{sec:feynman30-defects}, Table~\ref{tab:randomsplit},
Table~\ref{tab:pcasplit}. \textbf{Status: already incorporated} --- the
main manuscript's Table~\ref{tab:randomsplit} and Table~\ref{tab:pcasplit}
match this note's figures exactly (12/30 random-split, 13/30 PCA-split).
Reproduced here for the run-provenance disclosure, which is not otherwise
carried in the main text, and which directly resolves
\hyperref[supp:edit-C]{Edit Item C} in \S\ref{supp:tracker} below: the
Planck-blackbody near-miss/catastrophic divergence between splits is not
an unexplained anomaly in isolated tables, but part of a fully disclosed,
five-defect-catalogued run pair.
\end{appliesto}

\subsection*{Run provenance: which dates this section does, and does not, report on}

An investigation into a residual ``9/30 vs.\ 106/180'' denominator
question (\texttt{investigate\_section\_10\_7.py}) surfaced, in the
course of checking for the script's own intended canonical files, that
\textbf{at least four distinct PCA-split run dates} exist somewhere in
this repository's history. Table~\ref{supp:tab:provenance} discloses all
four explicitly, so that a reviewer encountering any of the other three
elsewhere in the repository does not reasonably mistake them for the
figures reported in this section.

\begin{table}[h]
\centering
\caption{Every PCA-split run date identified to date, and its status
with respect to Section~\ref{sec:feynman30}. Only the July~22/23 pair is
the source of the figures reported below.}
\label{supp:tab:provenance}
\begin{tabular}{lp{9.5cm}}
\toprule
Run date & Status \\
\midrule
June~4, 2026 & \textbf{Discarded, pre-fix.} Located under a legacy
\texttt{exp2/} directory rather than \texttt{exp2\_pca\_4060/}.
\texttt{HybridDiscoverySystem v50\_2 (tools)} scores 0/30, with all 30
equations flagged as environment/import failures (``not available''), not
accuracy failures --- the method did not execute at all in this run. Not
used anywhere in this section. \\
June~29, 2026 & \textbf{Not located.} Named as the canonical/current
batch by \texttt{investigate\_section\_10\_7.py}'s own configuration, but
absent from every file set supplied to date, at its canonical path and
under fallback search alike. Not used anywhere in this section; its
absence does not affect the figures below, which are independently
sourced from July~22/23. \\
June~30, 2026 & \textbf{Unverified, not used.} Source of
\texttt{exp2\_pca\_4060\_summary.json}'s pooled \texttt{n\_pass=142},
\texttt{n\_total=180} figure, which that file's own metadata pairs
explicitly with the already-withdrawn ``9/30 (random\_80\_20)'' claim.
No per-method recount was ever performed against this specific date's
protocol files (they were not among the dates checked), so 142/180
should continue to be treated as unconfirmed --- consistent with the
already-withdrawn 142/180 figure's known failure mode (arithmetic
pooling of records describing one run as if two) --- pending a direct
recount, which is out of scope for this update. Not used anywhere in
this section. \\
\textbf{July~22 / July~23, 2026} & \textbf{Canonical --- this is the run
pair Section~\ref{sec:feynman30} reports on.} July~22 (\texttt{protocol\_core\_extrap\_
20260722\_*.json}) is the random 80/20 split source; July~23
(\texttt{protocol\_core\_noiseless\_pca\_20260723\_*.json}) is the
PCA-directed 40/60 split source. Confirmed as canonical directly by the
corresponding author. \\
\bottomrule
\end{tabular}
\end{table}

\textbf{This distinction matters for citation hygiene, not just
correctness.} A reader who independently opens
\texttt{exp2\_pca\_4060\_summary.json} and sees 142/180, or who finds the
June~4 batch and sees \texttt{HybridDiscoverySystem} crash on every
equation, could otherwise reasonably conclude Section~\ref{sec:feynman30}
is internally inconsistent with the repository. It is not: those two
artifacts are from different, non-canonical dates, and neither has been
used to produce any figure below.

\subsection*{Headline figures (both gated on self-verification)}

\textbf{12/30 (40.0\%)} under the random 80/20 split (run dated
2026-07-22) and \textbf{13/30 (43.3\%)} under the PCA-directed 40/60
split (run dated 2026-07-23) --- the latter superseding an 18/30 figure
that had been carried unchecked through earlier drafts. Five confirmed
evaluator defects account for every null or blocked score across both
splits; one (a placeholder-token defect affecting 4 random-split
equations) remains architecturally unfixable by evaluator patching and is
disclosed as such below.

\subsection*{Full per-equation results}

\begin{longtable}{p{2.3cm}p{6.2cm}p{1.5cm}p{2.5cm}}
\caption{Random 80/20 split, run dated 2026-07-22, all 30 equations, after
patching three evaluator defects and gating on a 14/14 exact-reproduction
self-check. Bold: pass ($R^2 \geq 0.999999$). Matches main-text
Table~\ref{tab:randomsplit}.}
\label{supp:tab:randomsplit}\\
\toprule
Domain & Equation & \rtwo & Status \\
\midrule
\endfirsthead
\toprule
Domain & Equation & \rtwo & Status \\
\midrule
\endhead
Biology & Allometric scaling law & \textbf{1.0000} & pass \\
Biology & Logistic growth rate & \textbf{1.0000} & pass \\
Biology & Michaelis--Menten enzyme kinetics & \textbf{1.0000} & pass \\
Chemistry & Arrhenius rate constant & $-2.368$ & fail (finite) \\
Chemistry & Henderson--Hasselbalch buffer pH & \textbf{1.0000} & pass \\
Electrochemistry & Nernst equation & --- & unresolvable$^\ddagger$ \\
Electromagnetism & Clausius--Mossotti dielectric field & \textbf{1.0000} & pass \\
Electromagnetism & Dielectric polarisation & $-2.2\times10^{45}$ & catastrophic \\
Electromagnetism & Energy in a capacitor & $-1.240$ & fail (finite) \\
Electromagnetism & Lorentz force & $-\infty$ & catastrophic \\
Electromagnetism & Ohm's law & \textbf{1.0000} & pass \\
Electrostatics & Coulomb force, 1D simplified & $-0.066$ & fail (finite) \\
Electrostatics & Coulomb's law & $-0.253$ & fail (finite) \\
Magnetism & Curie's law & --- & unresolvable$^\ddagger$ \\
Mechanics & Kinetic energy (classical) & $-1.590$ & fail (finite) \\
Mechanics & Newton's gravitational force & 0.6634 & fail (finite) \\
Mechanics & Reduced mass, two-body system & \textbf{1.0000} & pass \\
Mechanics & Total mechanical energy & \textbf{0.9999998} & pass (recovered) \\
Optics & Double-slit interference intensity & \textbf{1.0000} & pass \\
Optics & Snell's law & \textbf{1.0000} & pass (recovered) \\
Probability & Gaussian probability density & 0.1340 & fail (finite) \\
Quantum & Bose--Einstein occupation number & --- & unresolvable$^\ddagger$ \\
Quantum & Fermi--Dirac occupation number & $-\infty$ & catastrophic \\
Quantum & Photon energy & $-\infty$ & catastrophic \\
Quantum & Rabi frequency & $-3.058$ & fail (finite) \\
Quantum & Zeeman energy & \textbf{1.0000} & pass \\
Thermodynamics & Fourier's law of heat conduction & $-0.2803$ & fail (finite) \\
Thermodynamics & Ideal gas law & \textbf{1.0000} & pass \\
Thermodynamics & Planck blackbody spectral radiance & 0.999994 & near-miss \\
Thermodynamics & Stefan--Boltzmann law & --- & unresolvable$^\ddagger$ \\
\midrule
\multicolumn{4}{l}{Pass: 12/30 (40.0\%). Near-miss: 1/30. Genuine fail: 13/30. Unresolvable: 4/30.} \\
\bottomrule
\end{longtable}

$\ddagger$ Placeholder-token defect (item~4 below): stored formula
references a variable name that was never a real input.

\bigskip

\begin{longtable}{p{2.3cm}p{6.2cm}p{1.5cm}p{2.5cm}}
\caption{PCA-directed 40/60 split, run dated 2026-07-23. The 17 rows
below were originally null and are individually re-derived in this
update, gated on a 13/13 exact-reproduction self-check against the 13
already-scored equations. \gap: the 13 already-scored equations' names
and individual \rtwo\ values were not captured at per-equation
granularity during this update's rescoring run (only the aggregate
10-pass / 3-fail split was recorded). Matches main-text
Table~\ref{tab:pcasplit}.}
\label{supp:tab:pcasplit}\\
\toprule
Equation & \rtwo & Status \\
\midrule
\endfirsthead
\toprule
Equation & \rtwo & Status \\
\midrule
\endhead
Allometric scaling law & \textbf{0.9999995} & pass (recovered) \\
Arrhenius rate constant & $-0.363$ & fail (finite) \\
Henderson--Hasselbalch buffer pH & 0.9998196 & near-miss \\
Nernst equation & 0.9999372 & near-miss \\
Dielectric polarisation & $-1.068$ & fail (finite) \\
Lorentz force & $-1.030$ & fail (finite) \\
Coulomb force, 1D simplified & $-0.863$ & fail (finite) \\
Coulomb's law & $-0.460$ & fail (finite) \\
Snell's law & \textbf{1.0000} & pass (recovered) \\
Gaussian probability density & $-1.192$ & fail (finite) \\
Photon energy & $-13.592$ & catastrophic \\
Bose--Einstein occupation number & $-1.051$ & fail (finite) \\
Rabi frequency & $-13.117$ & catastrophic \\
Planck blackbody spectral radiance & $-27.832$ & catastrophic \\
Fourier's law of heat conduction & $-0.084$ & fail (finite) \\
Stefan--Boltzmann law & $-0.889$ & fail (finite) \\
Ideal gas law & \textbf{0.9999999933} & pass (recovered) \\
\midrule
\multicolumn{3}{l}{Of these 17: 3 pass, 2 near-miss, 12 genuine fail.} \\
\bottomrule
\end{longtable}

\noindent\textbf{Remaining 13 equations (already-scored, not rescored
this round):} 10 pass, 3 fail, per the aggregate self-verification count.
\gap\ Individual equation names and \rtwo\ values for these 13 were not
printed during this update's Phase~1 run and are not reconstructed here
rather than guessed.

\subsection*{Defect catalog (cumulative, both splits)}

\begin{enumerate}
\item \textbf{\texttt{\textasciicircum} parsed as bitwise XOR, not
exponentiation.} Fixed by rewriting to \texttt{**} before evaluation.
Random split only.
\item \textbf{Unstripped bracketed annotation breaks the parser.} Fixed
by stripping a leading \texttt{[\ldots]} segment before evaluation.
Random split only.
\item \textbf{Missing \texttt{safe\_asin}/\texttt{safe\_acos}.} Fixed by
injecting clipped-inverse-trig definitions into the evaluator namespace.
Both splits.
\item \textbf{Placeholder-token variable names (e.g.\ \texttt{x0},
\texttt{x1}) baked into stored formulas.} \open\ Architecturally
distinct: no evaluator patch can recover a reference to a variable that
was never real. Affects 4 random-split equations (Nernst, Curie's law,
Bose--Einstein, Stefan--Boltzmann). Confirmed \emph{absent} from all 30
PCA-split formula strings by direct pattern search. Root cause (an
upstream formula-serialization step) not yet located.
\item \textbf{Missing \texttt{asin\_of\_sin}.} Newly identified this
round: 6 of 17 PCA-split null equations failed uniformly with
\texttt{NameError: name 'asin\_of\_sin' is not defined}. Confirmed via
five independent call sites in the production codebase as
$\texttt{asin\_of\_sin}(x)=\arcsin(\mathrm{clip}(\sin(x),-1,1))$ and
patched into the evaluator namespace. \open\ Patched only in the two
rescoring-time diagnostic evaluator copies; whether the \emph{production}
evaluator shares this gap has not been checked.
\end{enumerate}

\subsection*{Remaining open items}

\begin{enumerate}
\item \open\ Placeholder-token defect root cause (item~4 above) still
unfixed at the source; 4 random-split equations remain permanently
unevaluatable until the formula-serialization step is located.
\item \open\ Whether \texttt{asin\_of\_sin} is missing from the
production evaluator, not only the rescoring-time copies.
\item \open\ Near-miss threshold policy undecided: at a looser
convention, the random split reads 13/30 and the PCA split reads 15/30
instead of 12/30 and 13/30. This update uses the strict
$R^2\geq0.999999$ threshold throughout and reports near-misses
separately; this should be confirmed as the intended final convention.
\item \open\ June~30 run's 142/180 figure remains unverified (no direct
per-method recount performed against that specific date); it is
disclosed here as non-canonical rather than resolved.
\item \gap\ Per-equation identity and \rtwo\ values for the 13
already-scored PCA-split equations were not captured at the needed
granularity this round (Table~\ref{supp:tab:pcasplit}, note).
\end{enumerate}

\clearpage
% =======================================================================
\section{Reproducibility Note: Unpinned LLM Model String in the Nguyen-12 Pipeline}
\label{supp:nguyen12}
% =======================================================================

\begin{appliesto}
\S\ref{sec:nguyen12} (``Nguyen-12 Benchmark (Standard SR Suite)'').
\textbf{Status: not yet incorporated into main text} --- this note
resolves which model and sampling settings actually produced the
Nguyen-12 numbers (\texttt{claude-sonnet-5} at \texttt{temperature=0.25}
via \texttt{hypatia.py}, not \texttt{claude-sonnet-4-6} via
\texttt{LLMConfig} as previously assumed) and narrows, but does not
close, the 91.7\%-vs-100\% seed-123 discrepancy. Recommend folding the
``Revised recommendation'' paragraph below into
\S\ref{sec:nguyen12} or an adjacent reproducibility subsection, and
tracking the three action items as an explicit extension of issue
\texttt{FIX-N3}.
\end{appliesto}

While investigating a seed-reproducibility discrepancy in the Nguyen-12
extrapolation experiment (\S\ref{sec:nguyen12}, cf.\ FIX-N3 in the issue
tracker), we traced the actual language model invoked by the pipeline
across four source snapshots and found that it is not pinned to a dated
model release, which has direct consequences for whether currently
reported solve rates should be read as a reproduction of the originally
benchmarked figure. A fourth snapshot, \texttt{hypatia.py} --- the file
that turns out to sit directly on the Nguyen-12 benchmark's actual LLM
call path --- was obtained after the first three were traced and
resolves the model-identity question: see below.

\paragraph{What was checked.} Three artifacts spanning the project's
history were inspected for the model string actually compiled into the
symbolic-regression pipeline: (i) an executed notebook from an earlier
revision (\texttt{hybrid50v}); (ii) two Colab-exported scripts from a later
revision, one per seed (\texttt{hybrid50v\_02}); and (iii) the current
continuous-integration source, \texttt{symbolic\_engine.py}, as imported by
\texttt{run\_all.sh}. Each snapshot hardcodes a different model string as
the \texttt{LLMConfig} default (\texttt{claude-sonnet-4-6},
\texttt{claude-sonnet-4-5}, and \texttt{claude-sonnet-4-20250514},
respectively, the last of these appearing only in the CI configuration file
\texttt{repro.yaml} and an accompanying SDK-version comment). Critically,
none of the three code snapshots contains any mechanism to read an
environment variable or the \texttt{repro.yaml} \texttt{llm\_model} field at
runtime: \texttt{LLMConfig.model} is a plain dataclass default in each case,
and \texttt{hybrid\_system\_v50\_2.py} instantiates \texttt{LLMConfig}
without ever passing a \texttt{model} argument. The model that actually
produced any given run therefore depends entirely on which hardcoded source
revision happened to be present at execution time, not on the value
documented in the CI configuration.

\paragraph{Current live default, and a correction.} The presently shipping
source, \texttt{symbolic\_engine.py}, sets the default to
\texttt{claude-sonnet-4-6}. An earlier draft of this note hypothesized that
this string behaves as an undated, evergreen alias that Anthropic could
silently repoint to a different underlying model over time. Checking
Anthropic's published model-versioning documentation shows this is
incorrect: starting with the Claude 4.6 generation, dateless model IDs such
as \texttt{claude-sonnet-4-6} are themselves pinned, fixed snapshots, not
aliases --- Anthropic does not update the weights behind an existing model
ID, and a new version ships under a new ID entirely. Only pre-4.6 dateless
strings (for example \texttt{claude-sonnet-4-5}, which appears in the
\texttt{hybrid50v\_02} Colab exports rather than in the current CI source)
behave as convenience pointers that resolve to the latest dated snapshot
and can therefore drift. We retract the alias-drift explanation for the
currently-live code path and flag it explicitly so it is not repeated as
settled fact elsewhere in the paper or its supplementary material.

\paragraph{Resolved: the Nguyen-12 benchmark's actual live model is a
fourth, different string.} \texttt{hypatia.py} has now been obtained and
verified directly. \texttt{get\_llm\_prior()}'s own signature (lines
80--82) sets \texttt{n\_candidates: int = 5}, \texttt{temperature: float =
0.25}, and \texttt{model: str = "claude-sonnet-5"} --- a fourth distinct
hardcoded model string, different from all three traced previously
(\texttt{claude-sonnet-4-6} in \texttt{symbolic\_engine.py}'s
\texttt{LLMConfig}, \texttt{claude-sonnet-4-5} in the
\texttt{hybrid50v\_02}-era \texttt{LLMConfig} lineage, and
\texttt{claude-sonnet-4-20250514} in \texttt{repro.yaml}).
\texttt{exp3\_nguyen12\_hybrid50v\_02.py}'s call site (lines 455--459)
passes only \texttt{n\_candidates=8} and \texttt{verbose=False} --- no
\texttt{model=} or \texttt{temperature=} argument --- so the call falls
through to \texttt{hypatia.py}'s own defaults. The model that actually
produced the Nguyen-12 results is therefore \texttt{claude-sonnet-5} at
\texttt{temperature=0.25}, not \texttt{claude-sonnet-4-6} at
\texttt{temperature=0.3} as \texttt{symbolic\_engine.py}'s
\texttt{LLMConfig} would suggest; \texttt{n\_candidates=8} is a real,
confirmed override, but of \texttt{hypatia.py}'s own default of 5, not of
\texttt{LLMConfig}'s 3. There is a clear explanation for why
\texttt{claude-sonnet-4-6} was previously believed to be live:
\texttt{exp3\_nguyen12\_hybrid50v\_02.py:263} executes
\texttt{os.environ.setdefault('LLM\_MODEL', 'claude-sonnet-4-6')}, but
\texttt{hypatia.py} contains no reference to an \texttt{LLM\_MODEL}
environment variable anywhere --- the variable is set but nothing reads
it, the same ``looks wired in but isn't'' pattern already identified for
\texttt{repro.yaml}'s \texttt{llm\_model} field. \texttt{hypatia.py}'s
\texttt{client.messages.create()} call (lines 150--156) also passes no
seed parameter, independently confirming that no user-settable sampling
seed exists on this code path either. With \texttt{hypatia.py} resolved,
\texttt{run\_all.sh} --- to confirm which script CI actually invokes ---
is now the sole remaining unverified link in the chain.

\paragraph{Observed discrepancy (partially explained; timeline
narrowed).} The \texttt{hybrid50v\_02} script documents its own
validated target result in-line as 11/12 harmonic-rate solves (91.7\%),
10/12 pointwise-rate solves (83.3\%), 0/12 for the neural baseline, with a
Mann--Whitney comparison of $U{=}113$, $p{=}0.0097$. Recent reruns of the
same seed and code instead land at 12/12 (100\%). A cross-seed comparison
(seeds 99, 123, 777, 2024) run under confirmed identical full paper
parameters ($n_{\text{tasks}}{=}12$, $\text{niterations}{=}1000$,
$\text{populations}{=}30$, $\text{timeout}{=}1100\text{s}$) found 12/12
for every seed, which rules out both ``seed=123 was an unlucky draw'' and
``reduced sample size inflates the recovery rate'' as explanations. A
within-seed comparison at seed=123 across three timestamps found May~22
disagreeing with both July runs at the expression level on all 24
task/system entries (max $|r^2\ \text{diff}|\approx 8.46\times10^{-5}$),
while the two July runs (Jul~16, Jul~24) match each other exactly ---
something changed between May and mid-July and has been stable since. A
separately dated batch (filesystem-confirmed Apr~23, 2026, since none of
the five filenames matched the analysis script's own timestamp parser)
shows the pipeline was still running a reduced, single-task smoke-test
configuration at that point, and had moved to the full paper-parameter
configuration by May~22, 2026 --- narrowing the configuration-migration
window to roughly four weeks and flagging it as a plausible window for
the GCP/Kaggle~$\to$~GitHub~Actions infrastructure migration to have
occurred as well, since that migration remains the leading unexplained
variable once seed and sample size are ruled out. Sampling
nondeterminism at the LLM layer remains a contributing, independently
verified factor regardless of which code path (\texttt{LLMConfig} or
\texttt{hypatia.py}) is in play: \texttt{symbolic\_engine.py}'s own
module-level comment (lines 100--101) states outright that LLM
temperature sampling is server-side and cannot be seeded there, and its
\texttt{LLMConfig} default sets \texttt{temperature=0.3} (line 526) ---
a non-zero value with no exposed sampling seed. The Claude API is not
guaranteed to return bit-identical completions for identical prompts even
at low temperature, so if \texttt{seed} in this pipeline constrains only
the symbolic-regression search (PySR/NumPy) and not LLM sampling, two
runs with the ``same seed'' can still produce different LLM-proposed
candidate expressions, which PySR then refines into different final
fits. This is independently confirmed behavior of the traced
\texttt{LLMConfig} path; whether it also describes \texttt{hypatia.py}'s
behavior is unknown pending that file.

\paragraph{Corroborating billing evidence, re-read.} The API-key cost
dashboard still shows a key created shortly before the earlier, disagreeing
run with comparatively low cumulative cost, and a second key created about
seven weeks later that had, by the same last-used date, already accrued
roughly thirty times the cumulative cost of the first over a comparable
usage window. This remains useful evidence of \emph{something} changing
around that time --- additional retries, a higher candidate count, or a
configuration change --- but should now be read as pointing at a
configuration or usage-pattern change rather than at a model-version
change, since the model string itself is confirmed pinned.

\paragraph{Revised recommendation.} We no longer recommend ``pinning to a
dated snapshot'' as a fix for \texttt{symbolic\_engine.py}'s
\texttt{LLMConfig}, since \texttt{claude-sonnet-4-6} is already a fixed
snapshot there --- but that dataclass is now confirmed \emph{not} to be
what generated the Nguyen-12 numbers in the first place. The model that
matters, \texttt{claude-sonnet-5} in \texttt{hypatia.py}, is itself
already a dateless post-4.6-generation string and, by the same
documentation cited above, should likewise be a pinned snapshot rather
than a drifting alias, so model-identity drift remains an unlikely
explanation for the 91.7\%-vs-100\% gap on this corrected code path too.
The concrete action items are now: (1)~delete or actually wire up the
dead \texttt{LLM\_MODEL} environment variable in
\texttt{exp3\_nguyen12\_hybrid50v\_02.py}, since it currently misleads
readers of the code into believing \texttt{claude-sonnet-4-6} governs a
call it has no effect on; (2)~treat \texttt{hypatia.py}'s
\texttt{temperature=0.25} together with its unseeded
\texttt{client.messages.create()} call as the confirmed source of
run-to-run expression variability at fixed recovery rate, and record a
distribution over repeated seed=123 runs rather than a single 91.7\%
figure if that value is to be re-benchmarked; and (3)~obtain
\texttt{run\_all.sh}, which remains the sole unverified link --- it alone
can confirm whether CI invokes
\texttt{exp3\_nguyen12\_hybrid50v\_02.py} directly (as analyzed here) or
some other entry point that might call \texttt{get\_llm\_prior()} with
different arguments. Independently of the code-tracing question,
cross-seed and within-seed evidence from repeated
\texttt{seed\_agreement.py} runs (see above) already rules out seed
variance and reduced sample size as explanations for the
91.7\%-vs-100\% gap, and narrows the likely configuration-migration
window to Apr~23--May~22, 2026, which remains worth checking for an
infrastructure diff (e.g., the suspected
GCP/Kaggle~$\to$~GitHub~Actions migration) independent of the model
question, which is now closed. This does not change any number
currently printed in the paper text, which already states the
33.3\%/91.7\% split directly in prose. It does mean a future seed-123
result file should record \texttt{hypatia.py}'s sampling temperature and
candidate count alongside the harmonic- and strict-rate fields, not a
resolved model string from \texttt{symbolic\_engine.py}, so the
automated gate can check for the parameters that actually matter here.
Tracked as an extension of FIX-N3 pending the seed-123 result file and
\texttt{run\_all.sh}.

\clearpage
% =======================================================================
\section{CI Dashboard Status: Item 10.1}
\label{supp:item10-1}
% =======================================================================

\begin{appliesto}
\S\ref{sec:process_items} (``Outstanding Process/Credibility Items (Do Not
Cite Until Resolved)''), and Supplementary Material A \S10.1.
\textbf{Status: open, does not block submission} provided the CI
dashboard's ``0 open issues'' claim continues to go uncited elsewhere in
the paper. This section is the interim status note for that item.
\end{appliesto}

\noindent\textbf{Status: code fix confirmed; live re-run still pending.}
Not yet closeable --- do not cite the CI dashboard or its ``0 open issues''
claim until the checklist below is satisfied.

\subsection*{What was found}
Two notebooks were inspected directly:
\begin{itemize}
  \item \texttt{NB-06\_Code\_Quality\_Pipeline\_Integrity.ipynb} --- \textbf{0 of 10}
    code cells have ever executed (\texttt{execution\_count} is \texttt{null}
    throughout). This notebook has never actually run.
  \item \texttt{NB-04\_Numerical\_Consistency\_Checker.ipynb} --- execution
    stopped partway through (cells ran as \#1, \#3, \#4, \#5, then silently
    stopped). The cells that never ran include the notebook's own final
    \texttt{NB-04 STATUS} cell, which is the cell that declares
    \texttt{FIX-N1}, \texttt{FIX-N2}, and \texttt{FIX-C3} ``RESOLVED.''
    That declaration was never produced this cycle.
\end{itemize}

\noindent In \texttt{build\_report.py}, the dashboard's ``Open Issues''
table and its ``0 open issues'' verdict are computed from
\texttt{issue\_registry.json} via \texttt{load\_registry()} ---
a code path entirely independent of whether the notebooks above
existed or executed. A missing or never-run notebook only produced a
soft warning \texttt{<div>} in its own display section; it never
affected the findings table. So a stale or cached registry could
assert ``0 open issues'' while the notebooks meant to verify that
had not run at all.

\subsection*{Fix applied}
\texttt{build\_report.py} now checks every code cell's
\texttt{execution\_count} for each of the six notebooks (not just
whether the file exists), and \textbf{hard-fails the build}
(\texttt{sys.exit(1)}) if any notebook is missing or incomplete,
unless the run is an explicit fix-only pass (\texttt{NO\_NOTEBOOKS=true}).
Verified against the two real notebooks above (build correctly exits~1)
and against a synthetic fully-executed set (build correctly exits~0);
the \texttt{NO\_NOTEBOOKS=true} bypass was also confirmed to still work
as intended.

\subsection*{What is still open}
The code fix prevents a \emph{future} stale-green dashboard. It does not,
by itself, tell us what the six notebooks' findings actually are right now.
To close this item:
\begin{enumerate}
  \item Re-run \texttt{ci\_paper\_notebooks.yml} so all six notebooks
    execute to completion.
  \item Re-run \texttt{build\_report.py} (now patched) against that output.
    If it exits~0, the dashboard it produces is live and citable.
  \item Only then may the CI dashboard's conclusions --- including any
    ``N open issues'' figure --- be cited in the paper or supplementary
    material.
\end{enumerate}

\noindent Until that re-run happens, this item remains open per
Supplementary Material A, \S10.1, and per \S\ref{sec:process_items}
(``Outstanding Process/Credibility Items''). This does not block
submission provided the dashboard and its ``0 open issues'' claim
continue to go uncited elsewhere in the paper.

\clearpage
% =======================================================================
\section{Consolidated Manuscript-Edit Tracker}
\label{supp:tracker}
% =======================================================================

\begin{appliesto}
Cross-cutting: indexes every open or historical edit item against the
sections above and against the main manuscript. Confidence tiers: \textbf{A
--- mechanical} (the paper already has the correct number somewhere else
in the same file; safe to apply directly); \textbf{B --- needs judgment}
(text drafted, framing decision required); \textbf{C --- can't draft}
(flagged only; underlying fact not yet established).
\end{appliesto}

\subsection*{A. The \texttt{89.2\%} $\to$ \texttt{90.5\%} figure}

\textbf{Tier A, mechanical. Status: \closed\ applied in
\texttt{jmlr\_paper\_main\_3107919.tex}.} A prior
\texttt{grep} against
\texttt{jmlr\_paper\_main\_3007217.tex} confirmed \texttt{89.2\%} appeared
at seven locations (abstract footnote, note near
l.\,251, Contributions item~ii, Table~\ref{tab:main_results} caption and
cells, prose near the ablation discussion, and the Conclusion), while
Table~\ref{tab:difficulty} (l.\,1393) already reported the correct
\textbf{90.5\%}. That internal inconsistency has since been resolved: all
seven instances now read \textbf{90.5\%} in
\texttt{jmlr\_paper\_main\_3107919.tex}, matching the already-correct
Table~\ref{tab:difficulty} row, under the umbrella of
\S\ref{sec:hybrid-attribution-bug}, where the corrected \textbf{60.8\%}
figure lives alongside the uncorrected \texttt{90.5\%} figure it is
contrasted with. Two dependent derived figures were also corrected as
part of the same fix: the abstract's percentage-point gains (previously
$+27$pp/$+83.8$pp off the stale 89.2\% base, now $+28.3$pp/$+85.1$pp off
90.5\%). One instance was deliberately left unchanged: the
``previously published \texttt{89.2\%}/66-of-74 figure'' paragraph
reconciling a historical run against the recomputed 67/74, which
correctly cites 89.2\% as a historical value under investigation, not as
a current claim. No other number in these passages needed to
change --- the 60.8\% corrected figure, the 62.2\% LLM baseline, the
$-4.8$pp hard-tier figure, and the 22-masked-catastrophic count were
already correct and independently verified against the raw file.

\subsection*{B. Runtime / speedup section}

\textbf{Tier B originally, now \closed\ superseded.} This item is fully
resolved by \S\ref{supp:runtime} above and by main-text
\S\ref{sec:timing}: the $1.73\times$/$1.64\times$ speedup claim and the
73\% reduction claim are both withdrawn, and the complete
5-seed$\times$74-task$\times$2-split sweep shows hybrid $2.5$--$9.3\times$
\emph{slower} than the neural baseline on every seed and every split. Do
not re-apply any draft of this item that still asserts a positive
speedup, cites only a single seed, or treats the multi-seed sweep as
incomplete --- all three conditions were true at an earlier stage of this
audit and are no longer true.

\subsection*{C. Planck blackbody cross-split divergence}
\label{supp:edit-C}

\textbf{Tier A, additive --- Status: \closed\ figures now fully disclosed,
mechanism still \open.} Planck blackbody spectral radiance is a
near-miss under the random 80/20 split ($\Rsq=0.999994$,
Table~\ref{tab:randomsplit}) but catastrophic under the PCA-directed
40/60 split ($\Rsq=-27.832$, Table~\ref{tab:pcasplit}). Both figures are
now confirmed, provenance-disclosed, and defect-catalogued in
\S\ref{supp:feynman30} above (Section~10.7 update) rather than resting on
a single flagged footnote. We still do not have an explanation for
\emph{why} the two splits diverge this sharply on this one equation ---
one candidate mechanism worth checking is whether the PCA split's test
region falls in the numerically unstable Wien-tail limit of the Planck
function, per the original suggested footnote text --- but the
observation itself is no longer merely flagged; it is fully reproducible
from Table~\ref{supp:tab:randomsplit} and Table~\ref{supp:tab:pcasplit}
above.

\subsection*{D. PCA \texttt{random\_state} documentation}

\textbf{Tier C --- nothing to draft yet.} Not investigated in this audit.
No manuscript text is suggested because the underlying fact (whether the
PCA split is currently seeded/deterministic) hasn't been checked. If this
is to be closed out, the next step is to check
\texttt{sklearn.decomposition.PCA(\ldots)} instantiations in the
split-construction code for a fixed \texttt{random\_state}, then add one
sentence to the reproducibility/methods section stating the value used.

\subsection*{E. Second \texttt{exp2\_run.log}}

\textbf{Tier C --- nothing to draft yet.} Same situation: the
manuscript's own closing section already defers this to Supplementary
Material A \S10, which was not available for this audit. No manuscript
edit is suggested until that file is available to check.

\subsection*{F. Nguyen-12 model versioning (new this round)}

\textbf{Tier B --- \closed\ folded into
\texttt{jmlr\_paper\_main\_3107919.tex}.} A reproducibility-note paragraph
covering the \texttt{hypatia.py}/\texttt{claude-sonnet-5} finding, the dead
\texttt{LLM\_MODEL} environment variable, and the 91.7\%-vs-100\%
seed-123 discrepancy has been added to \S\ref{sec:nguyen12}, tracked
there as an explicit extension of \texttt{FIX-N3} pending recovery of
\texttt{run\_all.sh}, per the three action items recommended in
\S\ref{supp:nguyen12} above.

\subsection*{G. Table 6 / Core-15 ablation (already applied)}

\textbf{Tier A --- \closed.} See \S\ref{supp:table6} above. The main
manuscript's Table~\ref{tab:llm_ablation} already states it was rebuilt
from \texttt{\_merged.json}. No further action needed unless
Table~\ref{tab:llm_ablation} is re-edited from a source other than
\texttt{\_merged.json}, in which case it should be re-checked against the
\emph{Real} columns in \S\ref{supp:table6}.

\subsection*{H. CI dashboard, Item 10.1 (already applied)}

\textbf{Tier B --- \open, does not block submission.} See
\S\ref{supp:item10-1} above and \S\ref{sec:process_items} in the main
text. Remains open pending a completed re-run of
\texttt{ci\_paper\_notebooks.yml} and \texttt{build\_report.py}; the
``0 open issues'' claim must continue to go uncited elsewhere in the
paper until then.

\end{document}
